\documentclass[tikz,border=4pt]{standalone}
% Paleta y fuentes compartidas por el libro y las figuras (única fuente de verdad).
\RequirePackage{xcolor}
\definecolor{oxido}{HTML}{B7410E}
\definecolor{cobre}{HTML}{D27D2D}
\definecolor{carbon}{HTML}{2B2B2B}
\definecolor{azul}{HTML}{1F5F8B}
\definecolor{arena}{HTML}{F4EDE4}

\usepackage{fontspec}
\setsansfont{FiraSans}[Extension=.otf,UprightFont=*-Regular,BoldFont=*-SemiBold,ItalicFont=*-Italic,BoldItalicFont=*-SemiBoldItalic]
\setmonofont{FiraMono}[Extension=.otf,UprightFont=*-Regular,BoldFont=*-Medium,Scale=0.95]
\usetikzlibrary{positioning,arrows.meta,calc}
\begin{document}
\sffamily\small
\begin{tikzpicture}[>={Stealth[length=2.4mm]},
  cod/.style={anchor=west,font=\ttfamily\small,inner sep=1pt},
  estado/.style={draw=carbon,fill=arena,rounded corners=3pt,
    minimum width=5.4cm,minimum height=1.0cm,align=center,
    inner sep=4pt},
  nota/.style={font=\sffamily\footnotesize,text=carbon!75},
  flecha/.style={->,thick,carbon!80}]
  % Código de la función
  \fill[fondocodigo] (-0.2,0.4) rectangle (6.6,-2.8);
  \foreach \t [count=\i from 0] in {async fn pedir() -> u32 \{,
      \ \ \ \ let a = leer(1).await;,
      \ \ \ \ let b = leer(2).await;,
      \ \ \ \ a + b,
      \}} {\node[cod] at (0,-0.6*\i) {\t};}
  \node[nota,anchor=north west,align=left] at (-0.2,-3.1)
    {Modelo mental: el cuerpo equivale a un \texttt{enum}\\con un estado por cada punto de espera.};
  % Estados
  \node[estado] (s0) at (11.3,0.6)
    {\textbf{Inicio}\\[-1pt]\footnotesize aún no se ha sondeado};
  \node[estado,below=0.55cm of s0] (s1)
    {\textbf{Espera 1}\\[-1pt]\footnotesize guarda el futuro de \texttt{leer(1)}};
  \node[estado,below=0.55cm of s1] (s2)
    {\textbf{Espera 2}\\[-1pt]\footnotesize guarda \texttt{a} y el futuro de \texttt{leer(2)}};
  \node[estado,below=0.55cm of s2,fill=carbon,text=white] (s3)
    {\textbf{Terminado}\\[-1pt]\footnotesize devolvió \texttt{Ready(a + b)}};
  % Transiciones
  \draw[flecha] (s0) -- node[nota,right=3pt] {primer \texttt{poll}} (s1);
  \draw[flecha] (s1) -- node[nota,right=3pt] {\texttt{leer(1)} da \texttt{Ready(a)}} (s2);
  \draw[flecha] (s2) -- node[nota,right=3pt] {\texttt{leer(2)} da \texttt{Ready(b)}} (s3);
  % Bucles de Pending
  \foreach \s in {s1,s2} {
    \draw[->,thick,oxido] (\s.west) ++(0,0.22) arc[start angle=60,end angle=300,radius=0.26cm];
    \node[nota,text=oxido,anchor=east,align=right] at ($(\s.west)+(-0.65,0)$)
      {\texttt{Pending}:\\sigue aquí};
  }
\end{tikzpicture}
\end{document}
